\subsection{The arithmetic dictionary in 032, Sections 5--6}

The geometric target of the degree-one theta tensor is the Albanese
realization. The arithmetic extraction uses numerical characters of a
commuting algebra of algebraic Hecke endomorphisms; these are not the
global splitting characters $\mu_i$ or the torus inputs $\pi_i$.
Here are the exact formulas, with the conventions that determine them.

At a suitable split good prime with residue cardinality $q$, let
$\Phi$ denote the $q$-power Frobenius endomorphism of the reduced
Albanese. In 032, Proposition~5.5, p.~32, the integral endomorphisms
$D_1,D_2,D_3$ satisfy
\[
 \Phi^3-D_1\Phi^2+qD_2\Phi-q^3D_3=0.
\]
The Frobenius on crystalline degree-one cohomology is its pullback
$\varphi=\Phi^*$, with the same polynomial in $D_i^*$.
Replacing this by reciprocal Frobenius would invert the slopes and
would not be this formula. The valuation is normalized by
$\nu_\iota(q)=1$ in Section~5.2, p.~30.

Fix the finite theta input away from the active split place $w$.
Lemma~6.2, pp.~37--38, constructs $r\in E^\times$, $M>0$, and
$\omega\in\overline\Q^\times$ such that
\[
 r\bar r=1,\quad \operatorname{val}_w(r)=M,
 \quad\operatorname{val}_{\bar w}(r)=-M,
 \quad D_q^M=\omega,
 \qquad
 \omega^2=\zeta\prod_{\tau\in\Sigma}\tau(r)^{-3d(\tau)}.
\]
Here $D_q$ is the normalized radial dilation, the relation holds on
the image in theta classes, and $\zeta$ is a root of unity. The
square records the metaplectic ambiguity. The exponent $-3$ is
the active rank, not the dimension of the multiplicity line.

Proposition~6.3, pp.~39--40, gives the half-modulus projective
spherical parameter
\[
 \{q^{1/2},q^{-1/2},e\},\qquad
 e=e_0^{-1},\quad e_0^M=\omega,
 \qquad
 \nu_\iota(\sigma e)=-\frac32v(\sigma|_E).
\]
Every complex conjugate of $e$ has modulus one. A scalar-conjugate
active character has third parameter $e'=\pm\sigma e$.
For a full numerical character $\chi$, write its full parameter as
$b(q^{1/2},q^{-1/2},e')$. The central relation in Lemma~6.4
implies that $b^3e'$ is a root of unity. In equation~(6.9), p.~42,
\[
 \chi(D_1)=q e_1(\alpha),\qquad
 \chi(D_2)=q e_2(\alpha),\qquad
 \chi(D_3)=e_3(\alpha),
 \quad\alpha=b(q^{1/2},q^{-1/2},e').
\]
The $e_i$ are elementary symmetric polynomials. These are numerical
eigenvalues; the Hecke operators on their generalized spaces can
still have nilpotent parts.

Proposition~6.6, pp.~43--44, combines this polynomial, Weil weight,
filtered comparison, and weak admissibility. It gives the pure
Hodge index
\[
 t_\eta=\frac12-\frac13\nu_\iota(e')\in\{0,1\}.
\]
The generalized space is holomorphic for $t_\eta=1$ and
antiholomorphic for $t_\eta=0$. The proof uses polynomial
idempotents in extending algebraic Hecke endomorphisms, so the
projectors preserve the Hodge filtration. For each scalar
conjugate the filtration ranks are calculated over the actual
compositum of the definition field and the projector coefficient
field. Thus it does not assume independence of those fields or
semisimplicity of all ambient Hecke operators.

Theorem~5.1, p.~26, proved in Section~6.3, pp.~44--45, then applies
scalar descent to the whole selected generalized space. Formula
(6.13) reads
\[
 t_{\sigma\eta}=\frac{1+v(\sigma|_E)}2.
\]
It supplies the exact rational-Hodge maps from $U(v)$ into
$H^1(X,\Q)$; Corollary~6.7, p.~45, places a finite collection
on a common principal cover. Only then does the Albanese convert
the maps into algebraic maps to the prescribed CM abelian varieties.
This completes the original rank-three dictionary. A6 uses Liu's
all-rank component theorem to obtain the corresponding identification
in higher ball dimensions, without claiming that the displayed
cubic Frobenius calculation itself has already been generalized.

The relation to the definite machine is therefore at two levels:
the oscillator kernels and splitting rules are shared, whereas the
selected numerical Hecke eigenspaces, their integral projectors,
and the definite class-function normalization have to be matched.
Q9 makes the projector denominators explicit when their integral
matrices and annihilating polynomials have been supplied.
